% =====================================================================
%  FICHE 8 — THÈME 5 — NIVEAU FACILE — ÉNONCÉS
% =====================================================================
\documentclass[11pt,a4paper]{article}
\usepackage[utf8]{inputenc}
\usepackage[T1]{fontenc}
\usepackage{lmodern}
\usepackage[french]{babel}
\usepackage{amsmath,amssymb,mathtools}
\usepackage{icomma}
\usepackage{siunitx}
\sisetup{output-decimal-marker={,},per-mode=power,group-separator={\,},group-minimum-digits=5}
\usepackage[version=4]{mhchem}
\usepackage{xcolor}
\usepackage{tikz}
\usepackage[most]{tcolorbox}
\usepackage{enumitem}
\usepackage{array}
\usepackage{fancyhdr}
\usepackage[hidelinks]{hyperref}
\usetikzlibrary{arrows.meta,positioning,calc}
\usepackage[a4paper,margin=2.1cm,headheight=26pt,headsep=10pt]{geometry}
\usepackage{titlesec}

\definecolor{primary}{RGB}{150,98,162}
\definecolor{primarydark}{RGB}{92,52,110}

\tcbset{boxbase/.style={enhanced,breakable,boxrule=0.9pt,arc=2.5pt,left=3mm,right=3mm,top=2.3mm,bottom=2.3mm,fonttitle=\bfseries,coltitle=white,toptitle=0.6mm,bottomtitle=0.6mm}}
\newtcolorbox[auto counter]{exercice}[1][]{boxbase,colback=primary!4,colframe=primary,title={\textsc{Exercice}~\thetcbcounter\ \ifx\relax#1\relax\relax\else\textmd{--- \itshape #1}\fi}}

\pagestyle{fancy}
\fancyhf{}
\fancyhead[L]{\small\textcolor{primarydark}{\textbf{Fiche 8} — Thème 5 : Équilibre chimique et Le Chatelier}}
\fancyhead[R]{\small\textcolor{primarydark}{\textsc{Facile}}}
\fancyfoot[C]{\small\thepage}
\renewcommand{\headrulewidth}{0.8pt}
\renewcommand{\headrule}{\hbox to\headwidth{\color{primary}\leaders\hrule height \headrulewidth\hfill}}
\setlength{\parindent}{0pt}\setlength{\parskip}{3pt}

\begin{document}

\begin{center}
\begin{tikzpicture}
\node[rectangle,rounded corners=7pt,fill=primary,text=white,minimum width=16cm,
      minimum height=1.1cm,align=center,inner sep=6pt]
  {{\large\bfseries Thème 5 — Équilibre chimique et Le Chatelier}\\[1pt]{\small Niveau \textsc{Facile} — énoncés}};
\end{tikzpicture}
\end{center}

\begin{exercice}[écrire la constante d'équilibre]
Écrire l'expression de la constante d'équilibre $K$ pour chaque réaction (phase homogène).
\begin{enumerate}[label=\alph*),leftmargin=2em,topsep=2pt,itemsep=2pt]
  \item $\ce{N2}(g) + 3\,\ce{H2}(g) \rightleftharpoons 2\,\ce{NH3}(g)$ ;
  \item $\ce{PCl5}(g) \rightleftharpoons \ce{PCl3}(g) + \ce{Cl2}(g)$ ;
  \item $\ce{H2}(g) + \ce{I2}(g) \rightleftharpoons 2\,\ce{HI}(g)$.
\end{enumerate}
\end{exercice}

\begin{exercice}[équilibres hétérogènes : que garde-t-on ?]
Écrire $K$ en n'oubliant pas d'\textbf{exclure les solides et liquides purs}.
\begin{enumerate}[label=\alph*),leftmargin=2em,topsep=2pt,itemsep=2pt]
  \item $\ce{CaCO3}(s) \rightleftharpoons \ce{CaO}(s) + \ce{CO2}(g)$ ;
  \item $2\,\ce{A}(aq) + \ce{B}(s) \rightleftharpoons 2\,\ce{C}(g) + \ce{D}(aq)$ ;
  \item $\ce{FeO}(s) + \ce{CO}(g) \rightleftharpoons \ce{Fe}(s) + \ce{CO2}(g)$.
\end{enumerate}
\end{exercice}

\begin{exercice}[interpréter la valeur de $K$]
Pour trois réactions à une même température, on mesure les constantes d'équilibre suivantes. Indiquer, pour
chacune, si l'équilibre est déplacé vers les \textbf{produits}, vers les \textbf{réactifs}, ou s'il y a
coexistence comparable.
\begin{enumerate}[label=\alph*),leftmargin=2em,topsep=2pt,itemsep=2pt]
  \item $K = \num{1e5}$ ;
  \item $K = \num{2e-4}$ ;
  \item $K = 1{,}5$.
\end{enumerate}
\end{exercice}

\begin{exercice}[une perturbation à la fois]
La synthèse de l'ammoniac \;$\ce{N2}(g) + 3\,\ce{H2}(g) \rightleftharpoons 2\,\ce{NH3}(g)$\; est
\textbf{exothermique}. Le système est à l'équilibre. Indiquer le sens de déplacement (droite / gauche /
aucun) pour chaque action.
\begin{enumerate}[label=\alph*),leftmargin=2em,topsep=2pt,itemsep=2pt]
  \item On ajoute du $\ce{N2}$.
  \item On retire du $\ce{NH3}$.
  \item On augmente la pression (à température constante).
  \item On ajoute un catalyseur.
\end{enumerate}
\end{exercice}

\begin{exercice}[la température change $K$]
\begin{enumerate}[label=\alph*),leftmargin=2em,topsep=2pt,itemsep=2pt]
  \item Pour une réaction \textbf{exothermique}, comment varie $K$ si l'on \emph{augmente} $T$ ?
  \item Pour une réaction \textbf{endothermique}, comment varie $K$ si l'on \emph{augmente} $T$ ?
  \item Parmi ces trois facteurs, lequel modifie la \emph{valeur} de $K$ : la température, les concentrations
        initiales, ou l'ajout d'un catalyseur ?
\end{enumerate}
\end{exercice}

\end{document}
